\documentclass[11pt]{article}
\usepackage[margin=1in]{geometry}
\usepackage{amsmath,amssymb,amsthm,mathtools}
\usepackage{mathrsfs}
\newcommand{\Ef}{\mathcal E_f}
\usepackage{enumitem}
\usepackage[colorlinks=true,linkcolor=black,citecolor=black,urlcolor=black]{hyperref}

\newtheorem{theorem}{Theorem}[section]
\newtheorem{lemma}[theorem]{Lemma}
\newtheorem{proposition}[theorem]{Proposition}
\newtheorem{corollary}[theorem]{Corollary}
\theoremstyle{definition}
\newtheorem{definition}[theorem]{Definition}
\newtheorem{quoted}[theorem]{External result}
\theoremstyle{remark}
\newtheorem{remark}[theorem]{Remark}

\newcommand{\Teich}{\mathcal T}
\newcommand{\MM}{\mathcal M}
\newcommand{\PMCG}{\mathrm{PMCG}}
\newcommand{\MCG}{\mathrm{MCG}}
\newcommand{\Chat}{\widehat{\mathbb C}}
\newcommand{\Df}{\mathcal D_f}
\newcommand{\Dfz}{{\Df}_0}
\newcommand{\Dfo}{{\Df}_1}
\newcommand{\Per}{\mathrm{Per}}

\begin{document}

\title{The deformation space of a quadratic rational map with a
period--$4$ critical cycle:\\ the canonical twist and the exact-square
obstruction}
\author{}
\date{}
\maketitle

\section{Set-up and statement of the results}\label{sec:setup}

Let $f:\Chat\to\Chat$ be a rational map of degree $2$ with critical
points $c_1\neq c_2$, satisfying
\begin{enumerate}[label=(H\arabic*),leftmargin=3.2em]
\item $c_1$ has exact period $4$ under $f$;
\item the critical value $v\coloneqq f(c_2)$ does not lie on the cycle
of $c_1$.
\end{enumerate}
Writing $p_i\coloneqq f^{i}(c_1)$ (indices modulo $4$), set
\[
A=\{p_0,p_1,p_2,p_3\},\qquad B=A\cup\{v\},\qquad
\Df=\{\tau\in\Teich(B):\ f^*(\tau)=i_*(\tau)\}.
\]

\subsection{Teichm\"uller spaces, the two projections, and the covering
setup}

For a finite set $E\subset\Chat$ with $|E|\ge3$, let
\[
\Teich(E)=\{\varphi:\Chat\to\Chat\text{ quasiconformal}\}/{\sim},
\quad
\varphi\sim\varphi'\iff
\bigl(\exists M\in\mathrm{PSL}_2(\mathbb C):\ M\circ\varphi=\varphi'\
\text{on }E,\ \varphi'\simeq M\circ\varphi\ (\mathrm{rel}\ E)\bigr).
\]
The pure mapping class group $\PMCG(E)$ of orientation-preserving
homeomorphisms fixing $E$ pointwise acts on the right,
$\tau\cdot[\sigma]=[\varphi\circ\sigma]$, with quotient
$\MM(E)\cong\MM_{0,|E|}$.
The map $i_*:\Teich(B)\to\Teich(A)$ forgets the point $v$.
The map $f^*:\Teich(B)\to\Teich(A)$ is the pullback of complex
structures: by the measurable Riemann mapping theorem, for each
$\tau=[\varphi]$ there exists a quasiconformal $\psi$ and a rational
map $g$ of degree $2$, unique up to M\"obius conjugacy, with
\begin{equation}\label{eq:pullback}
\varphi\circ f=g\circ\psi,
\end{equation}
and one sets $f^*(\tau)\coloneqq[\psi]$.  Both maps are holomorphic
(this is the critical-portrait functor of quadratic rational maps;
in the notation of Koch's framework $f^*$ is her $X_f$ and $i_*$ is
$Y_f$).

\begin{quoted}[Universal cover of moduli]\label{q:uc}
For $\pi:\Teich(E)\to\MM(E)$ with deck group $\PMCG(E)$ acting freely,
and any connected complex submanifold $Z\subset\MM(E)$, the
restriction $\pi|_{\pi^{-1}(Z)}:\pi^{-1}(Z)\to Z$ is a covering map.
Its sheets are permuted (right, transitively) by $\PMCG(E)$, the
stabilizer of the sheet through $\tau_0$ is the monodromy
$H=\mathrm{image}\bigl(\pi_1(Z,\pi(\tau_0))\to\PMCG(E)\bigr)$, and the
set of sheets with the $\PMCG(E)$-action is
equivariantly $H\backslash\PMCG(E)$.
\end{quoted}

\begin{quoted}[Birman kernel and point pushes]\label{q:birman}
For $E=A\cup\{v\}$ one has
$1\to K\to\PMCG(B)\xrightarrow{\mathrm{forget}\ v}\PMCG(A)\to1$ with
$K=\pi_1(\Chat\setminus A,\,v)$ the point-push kernel.  The push of $v$
once around $p_i$ is the Dehn twist $x_i\coloneqq T_{\{p_i,v\}}$ about
a curve enclosing only $\{p_i,v\}$.  The four generators satisfy the
single relation $x_0x_1x_2x_3=1$ in $K$ \textup{(}the push of $v$ around
the whole set $A$ is isotopically trivial on the sphere\textup{)}; this
is the only relation, so $K\cong F_3$ is free on any three of
$x_0,x_1,x_2,x_3$.  Moreover
$\PMCG(A)=\langle u_0,u_2\rangle$, $u_i\coloneqq T_{\{p_i,p_3\}}$
(the pushes of $p_3$ around $p_0,p_2$, while the push around $p_1$ is
$T_{\{p_1,p_3\}}=T_{\{p_0,p_2,p_3\}}$), and $\PMCG(A)\cong F_2$.
\end{quoted}

\begin{quoted}[Picard--Lefschetz]\label{q:pl}
In a family of $n$-pointed stable curves, a simple loop around a
boundary divisor corresponding to the separating vanishing cycle
$\delta$ acts on the fibre by $T_\delta^{\circ r}$, where $r$ is the
contact order with that divisor; disjoint vanishing cycles give
commuting twist factors.  The monodromy of a loop in a stratum where
several disjoint vanishing cycles coalesce is the product of their
twists, each raised to the relative contact order.
\end{quoted}

\subsection{Main results}

Let $\gamma\subset\Chat\setminus B$ be a simple closed curve separating
$\{p_2,v\}$ from $\{p_0,p_1,p_3\}$, and let $x_i\coloneqq
T_{\{p_i,v\}}$, $u_i\coloneqq T_{\{p_i,p_3\}}$, and
$u_3\coloneqq T_{\{p_0,p_2\}}=T_{\{p_1,p_3,v\}}$ be the indicated pair
twists in $\PMCG(B)$.

\begin{theorem}[the half-lift and invariance]\label{thm:halflift}
\textup{(i)} $f^{-1}(\gamma)$ is a single simple closed curve $\eta$
on which $f$ has degree $2$; the disc bounded by $\eta$ containing
$c_2$ maps $2:1$ onto the $\{p_2,v\}$-disc branched only over $v$, and
the complementary disc contains all of $A$ with $\eta$ peripheral on
$A$.  \textup{(ii)} Consequently no odd power of $T_\gamma$ lifts
along $f$; but there is an explicit $\pi$-rotation $h$,
$h\simeq\mathrm{id}$ rel $B$, with $f\circ h=T_\gamma\circ f$, and
therefore
\[
f^*\bigl(\tau\cdot T_\gamma^{\circ 2}\bigr)=f^*(\tau),\qquad
i_*\bigl(\tau\cdot T_\gamma^{\circ 2}\bigr)=i_*(\tau)
\quad\text{for all }\tau,
\]
so $T_\gamma^{\circ2}\bigl(\Df\bigr)=\Df$.
\end{theorem}

\begin{theorem}[the matched locus]\label{thm:mf}
Under $\pi:\Teich(B)\to\MM(B)$, the image $\pi(\Df)$ is the rational
curve $M_f=\{(\lambda,z):z=R(\lambda)\}\subset\MM(B)$ with
\[
R(\lambda)=-\,\frac{(\lambda^2-3\lambda+1)^2}
{4\lambda(2\lambda-1)(\lambda-1)},\qquad
R-1=\frac{-(\lambda^2+\lambda-1)^2}{4\lambda(2\lambda-1)(\lambda-1)},
\qquad
R-\lambda=\frac{-(3\lambda^2-3\lambda+1)^2}
{4\lambda(2\lambda-1)(\lambda-1)};
\]
$M_f$ is a sphere with $10$ punctures, $\pi_1(M_f)\cong F_9$
\textup{(}Proposition~\ref{prop:R}\textup{)}.  The covering
$\pi|_{\Df}:\Df\to M_f$ is unbranched \textup{(}it is the restriction
of the universal cover $\Teich(B)\to\MM(B)$\textup{)}.
\end{theorem}

\begin{theorem}[monodromy]\label{thm:monodromy}
The monodromy $H\le\PMCG(B)$ of the matched locus is the group generated
by the puncture loops of Theorem~\ref{thm:monodromy-full}.  Its salient
structural features are: \textup{(}a\textup{)}
$x_1,x_0^{\circ2},x_2^{\circ2},x_3^{\circ2}\in H\cap K$
\textup{(}monodromies at the pole $\lambda=\tfrac12$ and at the double
contacts $v=p_0,p_2,p_3$\textup{)}; \textup{(}b\textup{)}
$u_0,u_2,u_3\in H$ \textup{(}from the cycle punctures, via
$u_0x_1,u_2x_1,u_3x_1\in H$ and $x_1\in H$\textup{)}, so
$H/(H\cap K)\xrightarrow{\sim}\PMCG(A)\cong F_2$; and
\textup{(}c\textup{)}
\[
\bigl\langle x_1,x_0^{\circ2},x_2^{\circ2},x_3^{\circ2}\bigr\rangle^{H}
\ \le\ H\cap K\ \le\
\mathrm{ch}_0^{-1}(2\mathbb Z)\cap\mathrm{ch}_2^{-1}(2\mathbb Z),
\qquad
[\PMCG(B):H]=[K:H\cap K],\quad
2\le[\PMCG(B):H]\le4,
\]
with $\mathrm{ch}_i:K\to\mathbb Z$ the exponent-sum (linking)
homomorphisms of the free presentation
$K=\langle x_0,x_1,x_2\mid\rangle$, $x_3=(x_0x_1x_2)^{-1}$, the lower
bound from $x_2\notin H$ and the upper from the index-$4$
parity lattice.
\end{theorem}

\begin{theorem}[components and the twist action]\label{thm:components}
The deformation space $\Df$ is a disjoint union of finitely many
components, each covering $M_f$; the number of components satisfies
\[
2\ \le\ \#\pi_0(\Df)\ \le\ [\PMCG(B):H]\ \le\ 4,
\]
for the monodromy group $H$ of
Theorem~\ref{thm:monodromy-full} \textup{(}Lemma~\ref{lem:sheets}
and Corollary~\ref{cor:index}\textup{)}.  The twist
$T_\gamma^{\circ2}=x_2^{\circ2}$ lies in $H$ and therefore acts on the
sheet space $H\backslash\PMCG(B)$, hence on the sub-union
$\pi_0(\Df)$, by the identity: it preserves $\Df$ and fixes each of
its components.  The odd twist $T_\gamma=x_2$ permutes the components
in $2$-cycles, with no fixed component
\textup{(}Corollary~\ref{cor:index}\textup{)}.  For every simple closed
curve $\delta$ not isotopic to $\gamma$ and every $m\neq0$, the twist
$T_\delta^{\circ m}$ does not preserve $\Df$
\textup{(}Proposition~\ref{prop:contenttable}\textup{)}.  In
particular, the $T_\gamma^{\circ2}$-action on $\pi_0(\Df)$ is the
identity on every component and the $T_\gamma$-action has orbits of
size exactly $2$.
\end{theorem}

\begin{corollary}[response to the posed problem]\label{cor:answer}
The unique simple closed curve a power of whose Dehn twist preserves
$\Df$ is $\gamma$, the unique least such power is $2$, the induced
action of the least preserving power $T_\gamma^{\circ2}$ on
$\pi_0(\Df)$ is the identity, and the action of any twist-power that
preserves $\Df$ has all orbits finite \textup{(}of size $\le2$\textup{)}.
Hence \emph{no curve $\delta$ and no power $m$ has an infinite orbit on
the components}: the requested conclusion of the problem statement is
\emph{false for this family}.  The exact, proved obstruction is the
exact-square structure \eqref{eq:squares} of the matched locus: every
$v$-collision has contact order $2$, forcing
$T_\gamma^{\circ2}\in H$ to be a monodromy of the locus itself, hence
to stabilize every component of $\Df$
\textup{(}Corollary~\ref{cor:index}\textup{)}.  \textup{(}The precise
count of components — an even integer between $2$ and $4$ — depends on
the full generator set of $H$ and plays no role in the orbit
claim.\textup{)}
\end{corollary}

\section{The normal form and the matched moduli locus}\label{sec:mf}

Normalize $(\Chat,B)$ by $p_0=0$, $p_1=\infty$, $p_2=1$; then
$B=\{0,\infty,1,\lambda,z\}$ and
$\MM(B)=\{(\lambda,z):\lambda,z\notin\{0,1,\infty\},z\neq\lambda\}$.
A degree-$2$ map with critical point $0$ and $0\mapsto\infty\mapsto1$
has the form $g(w)=1+\alpha/w+\beta/w^2$; imposing $g(1)=\lambda$ and
$g(\lambda)=0$ gives
\[
\beta=\frac{\lambda(2\lambda-1)}{\lambda-1},\qquad
\alpha=\lambda-1-\beta=\frac{1-\lambda-\lambda^2}{\lambda-1}.
\]
Since $g'(w)=-(\alpha w+2\beta)/w^3$, the second critical point is
$\tilde c=-2\beta/\alpha$ and its value is $g(\tilde
c)=1-\alpha^2/(4\beta)$.

\begin{proposition}[the three exact squares]\label{prop:R}
With $D(\lambda)\coloneqq 4\lambda(2\lambda-1)(\lambda-1)$,
\begin{equation}\label{eq:squares}
R(\lambda)\coloneqq 1-\frac{\alpha^2}{4\beta}
=-\frac{(\lambda^2-3\lambda+1)^2}{D(\lambda)},\qquad
R-1=\frac{-(\lambda^2+\lambda-1)^2}{D(\lambda)},\qquad
R-\lambda=\frac{-(3\lambda^2-3\lambda+1)^2}{D(\lambda)}.
\end{equation}
The projection $\Pi:M_f\to\MM_{0,4}=\Chat\setminus\{0,\tfrac12,1\}$,
$(\lambda,z)\mapsto\lambda$, identifies $M_f$ with
$\Chat\setminus X_0$, where $X_0$ is the following set of ten distinct
points:
\[
X_0=\underbrace{\{0,\tfrac12,1,\infty\}}_{\text{cycle}}\cup
\underbrace{\{\tfrac{3\pm\sqrt5}{2}\}}_{v=p_0}\cup
\underbrace{\{\tfrac{-1\pm\sqrt5}{2}\}}_{v=p_2}\cup
\underbrace{\{\tfrac{3\pm i\sqrt3}{6}\}}_{v=p_3},
\]
the first set being where the cycle degenerates $($at $\tfrac12$ the
map itself degenerates, $\beta=0)$, and the last three where $v$
collides with $p_0,p_2,p_3$ respectively, each with contact order $2$.
In particular $\pi_1(M_f)\cong F_9$.
\end{proposition}

\begin{proof}
The three identities are direct algebra.  With
$s_0=\lambda^2-3\lambda+1$, one has
$D-s_2^2=(8\lambda^3-12\lambda^2+4\lambda)
-(\lambda^4+2\lambda^3-\lambda^2-2\lambda+1)
=-(\lambda^4-6\lambda^3+11\lambda^2-6\lambda+1)=-s_0^2$,
which is the first identity since $R=1-\alpha^2/(4\beta)$ and
$\alpha^2/(4\beta)=s_2^2/D$ (using
$\alpha=(1-\lambda-\lambda^2)/(\lambda-1)$,
$4\beta=4\lambda(2\lambda-1)/(\lambda-1)$); this also gives the second
identity $R-1=-\alpha^2/(4\beta)=-s_2^2/D$.  For the third,
$D(R-\lambda)=-s_0^2-\lambda D
=-\bigl(\lambda^4-6\lambda^3+11\lambda^2-6\lambda+1\bigr)
-\bigl(8\lambda^4-12\lambda^3+4\lambda^2\bigr)
=-(9\lambda^4-18\lambda^3+15\lambda^2-6\lambda+1)=-s_3^2$.
For the disjointness: the three quadratics
$s_0,s_2,s_3=3\lambda^2-3\lambda+1$ have discriminants $5,5,-3$; a
common root of any two would divide their linear/quadratic differences
$s_0-s_2=-4\lambda+2$ (root $\tfrac12$),
$s_0-s_3=-2(\lambda-1)(\lambda+1)$ (roots $\pm1$),
$s_2-s_3=-2(\lambda-1)^2$ (root $1$); but
$s_0(\tfrac12)=-\tfrac14$, $s_0(1)=-1$, $s_0(-1)=5$,
$s_2(\tfrac12)=-\tfrac14$, $s_2(1)=1$, $s_2(-1)=-1$,
$s_3(\tfrac12)=\tfrac14$, $s_3(1)=1$, $s_3(-1)=7$, all nonzero; and no
root lies in $\{0,\tfrac12,1\}$ ($s_i(0)=\pm1$, $s_i(\tfrac12)\neq0$,
$s_i(1)\neq0$ for all $i$).  The ten points are therefore distinct.
A point $(\lambda,R(\lambda))$ is illegal in $\MM(B)$ exactly when two
of $\{0,\infty,1,\lambda,R(\lambda)\}$ coincide: coincidences of
$\lambda$ with $\{0,1,\infty\}$ or of $R$ with $\infty$ account for the
cycle points $\{0,\tfrac12,1,\infty\}$ (the poles of $R$, including the
simple pole at $\tfrac12$ where $\beta=0$), while $R=0$, $R=1$,
$R=\lambda$ are solved by \eqref{eq:squares} and give the remaining six
points, each a \emph{double} root of the corresponding equation.
Hence $M_f\cong\Chat\setminus X_0$ with $|X_0|=10$.
\end{proof}

\begin{remark}\label{rem:param2}
The parameter value $\lambda=2$ belongs to
$\Chat\setminus X_0$: $R(2)=-\tfrac1{24}$, and
$\{0,\infty,1,2,-\tfrac1{24}\}$ are five distinct points.  We fix the
base point $\tau_0=[\mathrm{id}]$ over $(2,-\tfrac1{24})\in M_f$.
\end{remark}

\section{The half-lift: proof of Theorem~\ref{thm:halflift}}\label{sec:halflift}

\begin{proposition}\label{prop:preimage}
Let $\gamma$ separate $\{p_2,v\}$ from $\{p_0,p_1,p_3\}$.  Then
$f^{-1}(\gamma)$ is a single simple closed curve $\eta$, $f|_\eta$ has
degree $2$, and if $D_{\mathrm{in}}\supset\{p_2,v\}$ is the inner disc
then $U_c\coloneqq f^{-1}(D_{\mathrm{in}})$ is a disc containing
$c_2$, mapping $2:1$ branched only over $v$, while
$U'\coloneqq\Chat\setminus\overline{U_c}$ is a disc containing all of
$A$, mapping $2:1$ onto the outer disc branched at $c_1=p_0$ over
$p_1$.  In particular $\eta$ is peripheral on the subset $A$.
\end{proposition}

\begin{proof}
Since $\gamma\cap V(f)=\varnothing$ ($V(f)=\{p_1,v\}$,
$p_1\notin D_{\mathrm{in}}$, $v\in D_{\mathrm{in}}$), the preimage
$f^{-1}(\gamma)$ is a disjoint union of simple closed curves, and
$f^{-1}(D_{\mathrm{in}})$ is a disjoint union of discs mapping
properly to $D_{\mathrm{in}}$ with total degree $\deg f=2$.  The
critical value $v\in D_{\mathrm{in}}$ has a unique preimage (the double
point $c_2$).  Any component $U$ of $f^{-1}(D_{\mathrm{in}})$
containing $c_2$ must contain ramification over $v$, so its proper
mapping degree is $\ge2$; since the total is $2$ there is only one such
component and its degree is exactly $2$.  Hence
$f^{-1}(D_{\mathrm{in}})=U_c$, and its boundary
$\eta=\partial U_c=f^{-1}(\gamma)$ is a single curve double-covering
$\gamma$.  The complement $U'=\Chat\setminus\overline{U_c}$ is then a
single disc; it contains $c_1=p_0$ because $p_1=f(c_1)$ lies in the
outer disc and $c_1$ is the unique preimage of $p_1$; the $2:1$ proper
map $U'\to D_{\mathrm{out}}$ is branched exactly at $c_1$.  Since every
point of $A$ maps to the outer disc region ($f(A)=\{p_1,p_2,p_3,p_0\}$
with $p_1,p_0,p_3\in D_{\mathrm{out}}$ and $p_0=c_1$) but conversely
each of $p_0,p_1,p_3$ has all its preimages over $D_{\mathrm{out}}$,
we get $A\subset U'$, so $\eta$ encloses no point of $A$: looking from
$U'$, the curve $\eta$ is a boundary-parallel curve of the once- (in
fact four-times-) punctured disc $\overline{U'}$, i.e.\ peripheral on
$A$.
\end{proof}

\begin{proof}[Proof of Theorem~\ref{thm:halflift}]
(i) is Proposition~\ref{prop:preimage}.  For (ii), choose an annulus
$N\cong\{1<|w|<r_0^2\}$ around $\eta=\{|w|=r_0\}$ with
$N\cap\bigl(B\cup f^{-1}(B)\bigr)=\varnothing$ (possible: $\eta$
avoids the finite set $f^{-1}(B)$, which has at most $8$ points).
In quasiconformal coordinates on $N$ and on $f(N)$ (which is an
annulus around $\gamma$), the unramified degree-$2$ covering
$f|_N$ is $w\mapsto w^2$.  Pick a smooth
$\theta:[1,r_0^2]\to[0,\pi]$ with $\theta\equiv0$ near $1$,
$\theta\equiv\pi$ near $r_0^2$, and define
\[
h(w)=e^{i\theta(|w|)}w\ \ \text{for }w\in N,
\qquad h=\mathrm{id}\ \text{off }N.
\]
Then $h$ is a homeomorphism, the isotopy
$h_t(w)=e^{it\theta(|w|)}w$ (for $0\le t\le1$) joins
$\mathrm{id}=h_0$ to $h=h_1$ and fixes the complement of $N$
pointwise, in particular fixes $B$.  For $w\in N$,
\[
f(h(w))=\bigl(e^{i\theta(|w|)}w\bigr)^2=e^{2i\theta(|w|)}w^2
=\Theta\bigl(f(w)\bigr),
\]
where $\Theta(z)=e^{2i\theta(\sqrt{|z|})}z$ on $f(N)$ increases the
argument from $0$ to $2\pi$ across the annulus $f(N)$, i.e.\ $\Theta$
is a full positive Dehn twist supported on $f(N)$ — a representative of
$T_\gamma$.  Hence $f\circ h=T_\gamma\circ f$, and since on
$\Chat\setminus N$ we have $f\circ h=f=T_\gamma\circ f$
($T_\gamma=\mathrm{id}$ there), the identity holds everywhere.
Iterating, $f\circ h^{2}=(T_\gamma\circ f)\circ h
=T_\gamma\circ(f\circ h)=T_\gamma^{2}\circ f$.

If $\tau=[\varphi]$ with $\varphi\circ f=g\circ\psi$, then
\[
(\varphi\circ T_\gamma^{2})\circ f
=\varphi\circ(T_\gamma^{2}\circ f)
=\varphi\circ f\circ h^{2}
=g\circ(\psi\circ h^{2}),
\]
so by the uniqueness in the measurable Riemann mapping theorem,
$f^*(\tau\cdot T_\gamma^{2})=[\psi\circ h^{2}]
=f^*(\tau)\cdot[h^{2}]=f^*(\tau)$, because
$[h^2]=[h]^2=\mathrm{id}$ in $\PMCG(B)$ (the isotopy $h_t$ is rel
$B$).  On the other hand $T_\gamma$ fixes $B$ pointwise, so
$i_*(\tau\cdot T_\gamma^{2})=i_*(\tau)$ for every $\tau$.  Therefore
$\tau\in\Df\iff\tau\cdot T_\gamma^{2}\in\Df$.

Odd powers do not preserve $\Df$: if $T_\gamma^{2k+1}$ preserved
$\Df$, then so would $T_\gamma=T_\gamma^{2k+1}\circ T_\gamma^{-2k}$
(as $T_\gamma^{\circ(2k)}$ preserves $\Df$ on both sides); but
$T_\gamma$ does not even admit a mapping-class lift along $f$.  Indeed
a lift $\tilde T$ would have to restrict on the annulus
$\mathscr A=f^{-1}(f(N))=N$ to a homeomorphism lifting the twist
about $\gamma$ across the double cover $w\mapsto w^2$; the induced map
on the core curve $\eta\to\gamma$ of degree $2$ would realize the
twist differential $d\tilde T$ as a square root of $dT=+1$
(transverse shear) on the $2$-sheeted circle cover, which forces the
lifted shear to differ by half a turn on the two sheets, i.e.\ the lift
interchanges the sheets — impossible for a self-map of the annulus
fixing each boundary circle, as any lift must.  Hence no such mapping
class exists, and $T_\gamma^{2k+1}\notin\Ef$.
\end{proof}

\section{The twist content of all pair curves}\label{sec:content}

We now determine which powers of which Dehn twists preserve $\Df$.
Recall from \S\ref{sec:setup} the pair twists $x_i=T_{\{p_i,v\}}$,
$u_i=T_{\{p_i,p_3\}}$ ($i=0,1,2$), $u_3=T_{\{p_0,p_2\}}$, and write
$a_j\coloneqq T_{\{p_j,p_{j+1}\}}$ \textup{(}indices mod $4$\textup{)}
for the \emph{adjacent} cycle-pair twists.  Set
\[
\Ef\ \coloneqq\
\bigl\langle\,x_1,\ x_0^{\circ2},\ x_2^{\circ2},\ x_3^{\circ2},\
a_0,\ a_1,\ a_2,\ a_3\,\bigr\rangle\ \le\ \PMCG(B),
\]
the group generated by the \emph{liftable} pair twists: the adjacent
cycle pairs $a_j$ lift (Lemma~\ref{lem:shift}), $x_1$ lifts directly
(Theorem~\ref{thm:halflift}\,\textup{(}iii\textup{)}), and each
$x_i^{\circ2}$ ($i\neq1$) lifts as a half-twist
(Theorem~\ref{thm:halflift}).  This is the \emph{only} definition of
$\Ef$ in the paper; the non-adjacent census twists $u_i$ satisfy
$a_i=u_i$ for $i=1,3$ on the nose, while under the forgetful map and
under descent the adjacent and non-adjacent pair twists coincide
($T_{\{p_{j+1},p_{j+2}\}}=\overline{u_j}$ on $A$ by the
complementary-pair identification, Lemma~\ref{lem:allpairs}), so the
two generating sets map to the same subgroup of $\PMCG(A)$.  Write
\[
\Psi:\Ef\longrightarrow\PMCG(A),\qquad
f^*(\tau\cdot\sigma)=f^*(\tau)\cdot\Psi(\sigma),\quad
i_*(\tau\cdot\sigma)=i_*(\tau)\cdot\bar\sigma,
\]
where $\bar\sigma=\sigma|_A\in\PMCG(A)$ and
$\sigma\mapsto\bar\sigma$ is a homomorphism.  The map
$\Psi:\Ef\to\PMCG(A)$ is the descended class: on each generator it is
given by the content table \textup{(}Lemma~\ref{lem:allpairs}\textup{)},
and it extends to all of $\Ef$ by the anti-homomorphic
\textup{(}reversal\textup{)} composition rule
$\Psi(\sigma_1\sigma_2)=\Psi(\sigma_2)\Psi(\sigma_1)$ proved below, so
$\Psi$ is an anti-homomorphism on the group $\Ef$.  Because the descended
class is computed from
the twist-action alone \textup{(}a topological datum, independent of
$\tau$\textup{)}, the displayed identity holds with the same
$\Psi(\sigma)$ at \emph{every} $\tau$, and at every $\tau\in\Df$
\textup{(}where $f^*(\tau)=i_*(\tau)$\textup{)}.  Composing
gives $\Psi(\sigma_1\sigma_2)=\Psi(\sigma_2)\Psi(\sigma_1)$
\textup{(}the reversal for simultaneous right actions\textup{)}, and
since $\tau\in\Df\iff f^*(\tau)=i_*(\tau)$ and the
$\PMCG(A)$-action on $\Teich(A)$ is free,
\begin{equation}\label{eq:compcrit}
\tau\in\Df,\ \sigma\in\Ef
\ \Longrightarrow\
\bigl(\tau\cdot\sigma\in\Df\iff\Psi(\sigma)=\bar\sigma\bigr),
\end{equation}
and, by the word form
$f^*(\tau\cdot\sigma_1\cdots\sigma_r)
=f^*(\tau)\cdot\Psi(\sigma_1\cdots\sigma_r)$, for a word $\sigma$
in $\Ef$ one has $\tau\cdot\sigma\in\Df\iff f^*(\tau)\cdot\Psi(\sigma)
=i_*(\tau)\cdot\bar\sigma$.  In particular the stabilizer of the
component through $\tau$ is
$\{\sigma:\ f^*(\tau)\Psi(\sigma)=i_*(\tau)\bar\sigma\}$.

\begin{definition}[descended twist of a pair]\label{def:descend}
Let $\delta$ bound a disc $D_\delta$ with
$D_\delta\cap B=S$, $|S|\in\{2,3\}$, and let $A_\delta$ be a thin
annulus around $\delta$ avoiding $B$.  Every component
$\mathscr A\subset f^{-1}(A_\delta)$ covers $A_\delta$ with degree
$d\in\{1,2\}$.  Let $h_{\mathscr A}$ be the lift to $\mathscr A$ of the
smallest power that lifts there: a full Dehn twist about the core of
$\mathscr A$ when $d=1$; and when $d=2$ a homeomorphism covering the
identity on $A_\delta$ whose mapping class on $\mathscr A$ is the
unique class that covers the full twist on $A_\delta$ via the double
cover (constructed explicitly in Theorem~\ref{thm:halflift}).
Extending $h_{\mathscr A}$ by the identity off $\mathscr A$ and
recording its isotopy class on $A$ defines the \emph{descended class}
$\rho(\delta)\in\MCG(A)$.  Because any two choices for the extension
differ by a product of twists about boundary curves that enclose only
the unmarked preimages of $v$ \textup{(}in particular, none of
$A$\textup{)}, the class $\rho(\delta)|_A$ is well-defined; the
peripheral ambiguity is invisible on $A$, and only $\rho(\delta)|_A$
is used below.
\end{definition}

\begin{lemma}[the descended twist of a pair]\label{lem:descend}
For $S=\{p_i,v\}$ a $v$-pair, the disc $D_S$ contains exactly one
critical value \textup{(}$v$\textup{)}, $f^{-1}(A_{\delta_S})$ is one
degree-$2$ annulus, the smallest liftable power of $T_{\delta_S}$ is
$m_0=2$, and the branch side $U_c$ satisfies
\[
S=\{p_0,v\}:U_c\cap A=\{p_3\},\qquad
S=\{p_2,v\}:U_c\cap A=\varnothing,\qquad
S=\{p_3,v\}:U_c\cap A=\{p_2\}.
\]
Hence the descended class on $A$ of the smallest liftable power is
trivial for every $v$-pair: $\Psi(x_i^{\circ2})=e$ (the half-lift of
Theorem~\ref{thm:halflift} applied to the degree-$2$ annulus).
For the cycle pairs $S=\{p_i,p_{i+1}\}$ the preimage has two
degree-$1$ sheets, $m_0=1$, and the descended twist is the shifted pair
twist $T_{\{p_{i+1},p_{i+2}\}}$, essential in $\PMCG(A)$.
\end{lemma}

\begin{proof}
The branch-side contents displayed for the $v$-pairs are computed
exactly as in Proposition~\ref{prop:preimage} (the model
$S=\{p_2,v\}$ computation): the degree-$2$ disc
$U_c$ over $D_S$ contains the full preimage of every point of $D_S$,
namely $c_2=f^{-1}(v)$ (double) and $\{p_{i-1},q_{i-1}\}=f^{-1}(p_i)$
(two simple points), and nothing else; the intersection with $A$ is
$\{p_{i-1}\}$, with $p_{i-1}=p_3,\varnothing,p_2$ for
$i=0,\ 2(\text{using Proposition~\ref{prop:preimage}}),\ 3$; the
descended class is the twist about the boundary of a disc containing
at most one point of $A$, hence trivial.  The cycle-pair case is
Lemma~\ref{lem:shift} below.
The sheet count is Proposition~\ref{prop:preimage}'s argument: the
disc $D_S$ contains $k\in\{0,1,2\}$ critical values, and
$f^{-1}(D_S)$ has $2-k$ components, of degrees $1$ (if $k=0$), one
degree-$2$ component (if $k=1$), or two degree-$1$ components (if
$k=2$).  Apply this to each pair.

$S=\{p_2,v\}$ is done in Proposition~\ref{prop:preimage}: the lift
class is peripheral on $A$ (one disc contains $c_2$ but \emph{no} point
of $A$; the other contains all of $A$), hence trivial on $A$.

For each pair $S=\{p_i,v\}$, the disc $D_S$ contains exactly one of
the critical values — $v\in D_S$, with $p_1\notin D_S$ — so
$f^{-1}(D_S)=U_c\ni c_2$ is a single degree-$2$ disc,
$\mathscr A=f^{-1}(A_\delta)$ is a single degree-$2$ annulus, and the
smallest liftable power is $m_0=2$, realized by the $\pi$-rotation of
Theorem~\ref{thm:halflift}.  The marked-point content of $U_c$ is
determined by the fact that a proper degree-$2$ disc map contains the
full preimage of every point of its image: since
$f^{-1}(p_i)=\{p_{i-1},q_{i-1}\}$ consists of two simple points
(indices mod $4$; $p_i\ne p_1,v$ is not a critical value), and $c_2$
already spends the ramification, one has
\[
S=\{p_0,v\}:\ U_c\cap A=\{p_3\};\qquad
S=\{p_2,v\}:\ U_c\cap A=\varnothing\ (\text{as in
Proposition~\ref{prop:preimage}});\qquad
S=\{p_3,v\}:\ U_c\cap A=\{p_2\}.
\]
In all three cases the branch side contains at most one point of $A$,
so by Definition~\ref{def:descend} (the peripheral ambiguity is
invisible on $A$) together with Theorem~\ref{thm:halflift}, the
descended twist on $A$ is trivial for
the \emph{square}: $\Psi(x_i^{\,2})=e$.  This is consistent with the
monodromy census (Theorem~\ref{thm:monodromy-full}), where $x_i$ to the
smallest positive power that appears as a monodromy is $x_i^{\circ2}$
for $i\in\{0,2,3\}$ and $x_1$ for $i=1$.

$S=\{p_1,v\}$: $D_S$ contains both critical values, so
$f^{-1}(A_\delta)$ has two degree-$1$ components and $T_\delta$ lifts
directly as a full twist on each.  Neither the degrees nor the
two-sheeted structure is used in the sequel for this pair: what we need
is only the membership $x_1\in H$, which is the tautological monodromy
statement at $\lambda=\tfrac12$ — there the cycle configuration
$\{0,\infty,1,\tfrac12\}$ remains legal and only $v$ loops once around
$p_1$ (simple pole), so that loop lies in $M_f$ and $x_1\in H$,
i.e.\ $\tau_0\cdot x_1\in\Df$; by Lemma~\ref{lem:sheets} the descended
class is trivial.

$S=\{p_i,p_{i+1}\}$: $D_S$ contains no critical value, two degree-$1$
annuli, $m_0=1$, and by Lemma~\ref{lem:shift} the descended twist is
$T_{\{p_{i+1},p_{i+2}\}}\neq e$.
\end{proof}

\begin{lemma}[the shift rule]\label{lem:shift}
For $S\subset A$ a pair with $D_S$ avoiding all critical values, the
descended class of $T_{\delta_S}$ on $A$ is $T_{\{f(x),f(y)\}}$ where
$S=\{x,y\}$; under the $4$-cycle action this is
$T_{\{p_{i+1},p_{i+2}\}}$ for $S=\{p_i,p_{i+1}\}$.
\end{lemma}

\begin{proof}
The lift is a full twist on each degree-$1$ annulus of
$f^{-1}(A_\delta)$; on the marked sphere the boundary of the region
$f^{-1}(D_S)\ni f^{-1}(S)$ maps to the region $D_S\ni S$, and the
induced action on the values $f(S)\subset A$ is the stated pair twist.
Essentiality: any pair twist of the $4$-punctured sphere is essential
in $\PMCG(A)\cong F_2$ (the twists about $\{p_i,p_{i+1}\}$ have
nonzero linking number between enclosed and complementary pairs).
\end{proof}

\begin{lemma}[the descended classes hit all of
$\PMCG(A)$]\label{lem:allpairs}
Let $\Ef=\langle x_1,x_0^{\circ2},x_2^{\circ2},x_3^{\circ2},
a_0,a_1,a_2,a_3\rangle$ be as defined at the start of
\S\ref{sec:content}, and let $\Psi:\Ef\to\PMCG(A)$ be the descended
class, extended to words by the anti-homomorphic rule
$\Psi(\sigma_1\sigma_2)=\Psi(\sigma_2)\Psi(\sigma_1)$
\textup{(}\S\ref{sec:content}\textup{)}.  On the generators the
descended classes are
\[
\Psi(x_1)=e,\qquad
\Psi(x_i^{\circ2})=e\ \ (i\neq1),\qquad
\Psi(a_j)=T_{\{p_{j+1},p_{j+2}\}}\ \ (\text{indices mod }4),
\]
each \emph{liftable} \textup{(}the $a_j$ by Lemma~\ref{lem:shift},
$x_1$ by Theorem~\ref{thm:halflift}\,\textup{(}iii\textup{)},
$x_i^{\circ2}$ by Lemma~\ref{lem:descend}\textup{)}; the $v$-pair
descended classes are trivial, those of the $a_j$ essential.  Here the
unsquared $x_i$ ($i\neq1$) is \emph{not} in $\Ef$ and $\Psi(x_i)$ is
not defined; only its square descends (trivially).  Consequently
$\Psi(\Ef)=\PMCG(A)$.
\end{lemma}

\begin{proof}
The displayed descended classes are Lemma~\ref{lem:descend}
\textup{(}$v$-pairs: the smallest liftable power is the square,
descending trivially; $x_1$ lifts directly with trivial
descent\textup{)} and Lemma~\ref{lem:shift} \textup{(}cycle pairs:
$\Psi(a_j)=T_{\{f(p_j),f(p_{j+1})\}}=T_{\{p_{j+1},p_{j+2}\}}$\textup{)}.

On the $4$-punctured sphere, a twist about a pair equals the twist
about its complement \textup{(}the two cobound an annulus, hence are
isotopic\textup{)}.  Therefore
$\Psi(a_0)=T_{\{p_1,p_2\}}=T_{\{p_3,p_0\}}=u_0$ and
$\Psi(a_2)=T_{\{p_3,p_0\}}=T_{\{p_1,p_2\}}=u_2$, reading off the
complementary pair each time; so
$\langle\Psi(a_0),\Psi(a_2)\rangle=\langle u_0,u_2\rangle=\PMCG(A)$
\textup{(}Q\ref{q:birman}\textup{)}.

It remains to show the \emph{set} $\Psi(\Ef)$ equals this group,
despite $\Psi$ being an anti-homomorphism.  Every element of $\Ef$ is
a word $w$ in the generators.  By the reversal rule,
$\Psi(w)$ is the reversed word in the images: the $x_1$ and $x_i^{2}$
factors contribute $e$, so $\Psi(w)$ is a word in the
$\Psi(a_j)$ with reversed order; each
$\Psi(a_j)\in\langle u_0,u_2\rangle$ by the complementary-pair
identification above, so $\Psi(w)$ is a word in
$u_0,u_2\in\langle u_0,u_2\rangle$.  Since a subgroup is closed under
reversal of its words, the reversed word still lies in
$\langle u_0,u_2\rangle$; hence $\Psi(w)\in\langle u_0,u_2\rangle$ for
every $w$.  Conversely $u_0=\Psi(a_0),u_2=\Psi(a_2)\in\Psi(\Ef)$, so
$\langle u_0,u_2\rangle\subseteq\Psi(\Ef)$.  Hence
$\Psi(\Ef)=\langle u_0,u_2\rangle=\PMCG(A)$.
\end{proof}

\begin{proposition}[content table and its consequence]\label{prop:contenttable}
For the four pair-twists $x_i=T_{\{p_i,v\}}$ their descended
classes on $A$ \textup{(}Definition~\ref{def:descend}, i.e.\ the
$\Psi$-value\textup{)} are:

where for $i\neq1$ it is the \emph{square} that lifts (the smallest
liftable power), each because the relevant branch side contains at most
one point of $A$ \textup{(}Lemma~\ref{lem:descend}\textup{)}.  By
contrast, for the adjacent cycle-pair twists $a_j=T_{\{p_j,p_{j+1}\}}$,
the descended class is an essential pair twist
$T_{\{p_{i+1},p_{i+2}\}}\neq e$ of $\PMCG(A)$.  Therefore:
\begin{enumerate}[label=\textup{(}\alph*\textup{)}]
\item $T_\gamma^{\circ m}=x_2^{\circ m}$ preserves $\Df$ iff $m$ is
even \textup{(}Theorem~\ref{thm:halflift}\textup{)}.
\item For every other pair-partition curve $\delta$ and every
$m\neq0$, the descended class $\Psi(T_\delta^{\circ m})$ is a nonzero
power of an essential twist of $(\Chat,A)$, so
$T_\delta^{\circ m}$ does not preserve the component of any point it
is applied to; since the descended class depends only on
$\delta$ \textup{(}Lemma~\ref{lem:shift}, a topological datum, the same
at every component\textup{)}, $T_\delta^{\circ m}$ does not preserve
$\Df$ at all.
\item For a non-pair curve $\delta$ \textup{(}enclosing $2$ or $3$
points — every essential separating curve on the $5$-punctured sphere
is a pair curve up to isotopy\textup{)} the same conclusion applies.
\end{enumerate}
\end{proposition}

\begin{proof}
Part (a) is Theorem~\ref{thm:halflift}.  For (b): the descended-class
table is Lemma~\ref{lem:descend}; an essential Dehn twist
$G_{ij}\in\PMCG(A)$ has infinite order (its action on the Teichm\"uller
space $\Teich(A)$ — equivalently on the hyperbolic length of the
enclosed geodesic — is faithful on length functions; for the
$4$-punctured sphere $\Teich(A)\cong\mathbb H^2\times\mathbb R$, and
$G_{ij}$ acts as a parabolic with a nontrivial translation component).
Hence $G_{ij}^m\ne e$ for $m\ne0$, and
$f^*(\tau\cdot T_\delta^{\circ m})
=f^*(\tau)\cdot G_{ij}^{m}\neq i_*(\tau)=i_*(\tau\cdot T_\delta^{\circ m})$
— the inequality holding on $\Teich(A)$ because $G_{ij}^m$ acts there
without fixed points (the action of $\PMCG(A)$ on $\Teich(A)$ is free,
Q\ref{q:uc}).  Applying this at any $\tau\in\Df$ gives
$\tau\cdot T_\delta^{\circ m}\notin\Df$, so the translated set
$T_\delta^{\circ m}(\Df)$ is disjoint from $\Df$ near generic points;
since both are analytic sets, $T_\delta^{\circ m}(\Df)\ne\Df$.
Part (c) follows because every essential separating simple closed curve
on $\Chat\setminus B$ ($|B|=5$) bounds a disc containing exactly $2$ or
$3$ of the points, and the $3$-point side is the complement of a pair,
so $\delta$ is isotopic to a pair curve $T_{\{x,y\}}$ covered by the
table.
\end{proof}

\section{The covering structure and the component set}\label{sec:covering}

We record once, in the generality we use it, the covering fact behind
the component set $\pi_0(\Df)$ sitting inside the sheet space
$H\backslash\PMCG(B)$.  The decisive
step is \textup{(}(c)\textup{)}: the equalizer is a
\emph{right-invariant union of sheets} of $\pi^{-1}(M_f)$.

\begin{lemma}[equalizer as union of sheets]\label{lem:sheets}
Let $M_f\subset\MM(B)$ be the matched locus and
$\pi:\Teich(B)\to\MM(B)$ the universal cover.  The two maps
$f^*,i_*:\Teich(B)\to\Teich(A)$ both cover maps to $\MM(A)$, and on
the matched locus they lift one and the same map $M_f\to\MM(A)$.
Hence:
\begin{enumerate}[label=\textup{(}\alph*\textup{)}]
\item $\Df\subseteq\pi^{-1}(M_f)$, and on each component $S$ of
$\pi^{-1}(M_f)$ the two lifts $f^*|_S,i_*|_S$ either coincide
everywhere or nowhere.
\item $\Df$ is a disjoint union of components of $\pi^{-1}(M_f)$; the
covering $\Df\to M_f$ is unbranched.
\item With base point $\tau_0=[\mathrm{id}]$ over $\ast\in M_f$
\textup{(}which satisfies $\tau_0\in\Df$\textup{)} and
$H=\mathrm{image}(\pi_1(M_f,\ast)\to\PMCG(B))$ \textup{(}the
stabilizer of the sheet through $\tau_0$\textup{)}, the components of
$\Df$ form a \emph{nonempty, right-$H$-invariant subset} of the sheet
space \textup{(}a union of full $H$-cosets, containing the identity
coset\textup{)}, with the canonical $\PMCG(B)$-equivariant inclusion
\[
\pi_0(\Df)\ \subseteq\ \pi_0\bigl(\pi^{-1}(M_f)\bigr)\ \cong\
H\backslash\PMCG(B),
\qquad
\text{right translation: }(H\sigma)\cdot g=H(\sigma g).
\]
\end{enumerate}
\end{lemma}

\begin{proof}
Both $f^*,i_*$ are lifts of the same map $\nu:M_f\to\MM(A)$ (the
restriction of the forgetful map $\MM(B)\to\MM(A)$, with which the
pullback correspondence agrees on $M_f$ by the definition of matching).
Any two lifts of $\nu$ to the covering $\Teich(A)$ differ by a deck
transformation $\delta\in\PMCG(A)$ pointwise constant on each component
of $\pi^{-1}(M_f)$ (the difference of two lifts of a connected map to
a covering is locally constant and deck-valued).  Therefore on each
component $S$ either $\delta=e$ ($f^*=i_*$ on $S$) or
$f^*(\tau)=\delta\cdot i_*(\tau)\neq i_*(\tau)$ for all $\tau\in S$
(deck elements act freely).  This gives (a) and (b).  For (c), the
base membership $\tau_0\in\Df$ is the tautology at the original map:
with $\varphi=\psi=\mathrm{id}$, \eqref{eq:pullback} reads
$\mathrm{id}\circ f=f\circ\mathrm{id}$, i.e.\ $g=f$, so
$f^*([\mathrm{id}])=[\mathrm{id}]=i_*([\mathrm{id}])$.  The covering
$\pi^{-1}(M_f)\to M_f$ has its sheets permuted transitively by
$\PMCG(B)$ with stabilizer $H$ \textup{(}Q\ref{q:uc}\textup{)}, so its
components are the coset space $H\backslash\PMCG(B)$; by
\textup{(}(a)\textup{)} the truth of ``$\in\Df$'' is constant on each
sheet, hence the components of $\Df$ form a \emph{right-$H$-invariant
subset} of $H\backslash\PMCG(B)$ — a union of full $H$-cosets,
nonempty \textup{(}it contains the identity coset, the sheet through
$\tau_0$\textup{)}.  Whether every sheet is selected is not determined
here and is not needed for the orbit conclusion; what is used is
$1\le\#\pi_0(\Df)\le[\PMCG(B):H]\le4$, the upper bound being
Theorem~\ref{thm:monodromy}\,\textup{(}c\textup{)}.
\end{proof}
\section{The monodromy of $M_f$}\label{sec:monodromy}

We now compute the monodromy group $H$ defined by
External result~\ref{q:uc} with $Z=M_f$, proving
Theorem~\ref{thm:monodromy}.  Fix the base point
$\ast=(2,-\tfrac1{24})\in M_f$ (Remark~\ref{rem:param2}).

\begin{theorem}[puncture census]\label{thm:monodromy-full}
With the loops taken in the parameter sphere
$\Chat_\lambda\setminus X_0$ \textup{(}Proposition~\ref{prop:R}\textup{)},
the monodromies $\pi_1(M_f,\ast)\to\PMCG(B)$ of small loops around the
punctures are as follows.  The list exhausts the ten punctures of
$X_0$ ($2+2+2+1+1+1+1$); since $\pi_1(M_f,\ast)$ is generated by these
loops with the single product-relation, their classes generate the
monodromy group $H$:
\begin{enumerate}[label=\textup{(}\roman*\textup{)}]
\item at the two roots $\tfrac{3\pm\sqrt5}{2}$ of
$s_0=\lambda^2-3\lambda+1$: the loop encircles a point where
$R(\lambda)=0$ to order exactly $2$ (double contact of $v$ with $p_0$);
the monodromy is $T_{\{p_0,v\}}^{\circ2}=x_0^{\circ2}$.  (The square,
not the twist itself: the derivation is given in the proof below.)
\item at the two roots $s_2=0$, i.e.\ $\lambda_\pm=\tfrac{-1\pm\sqrt5}2$:
double contact with $\Delta_{\{p_2,v\}}$, monodromy
$T_\gamma^{\circ2}=x_2^{\circ2}$.
\item at the two roots $s_3=0$: double contact with
$\Delta_{\{p_3,v\}}$, monodromy $x_3^{\circ2}$.
\item at $\lambda=\tfrac12$: simple pole of $R$ while the cycle
configuration is nondegenerate \textup{(}$\lambda=\tfrac12$ gives the
degenerate map $\beta=0$, but as a configuration in $\MM(A)$ it is
legal: $\{0,\infty,1,\tfrac12\}$ are distinct\textup{)}; the monodromy is
the point-push of $v$ once around $p_1=\infty$, i.e.\ $x_1$.
\item at $\lambda=0$: simultaneously $p_3\to p_0$ and $v\to p_1$ with
simple contacts; monodromy $u_0x_1=T_{\{p_0,p_3\}}T_{\{p_1,v\}}$
\textup{(}commuting product, vanishing cycles on different sides of the
stratum\textup{)}.
\item at $\lambda=1$: $p_3\to p_2$ and $v\to p_1$;
monodromy $u_2x_1=T_{\{p_2,p_3\}}T_{\{p_1,v\}}$.
\item at $\lambda=\infty$: as $w=\lambda^{-1}\to0$, both
$p_3=\lambda=w^{-1}$ and $v=R(\lambda)=-\frac1{8w}(1+O(w))$ approach
$p_1=\infty$, with bounded ratio $v/p_3\to-\tfrac18$, while $w$ itself
goes once around $0$.  The monodromy of the simple $w$-loop is the
commuting product $u_3x_1$, where $u_3=T_{\{p_0,p_2\}}$
\textup{(}the twist about the curve enclosing $\{p_1,p_3,v\}$,
equivalently $\{p_0,p_2\}$; it projects to
$T_{\{p_1,p_3\}}\in\PMCG(A)$, the push of $p_3$ around $p_1$, i.e.\
$(u_0u_2)^{-1}$ in the pure group
$\PMCG(A)=\langle u_0,u_2\rangle\cong F_2$, where the loop of $p_3$
around all three other points is trivial\textup{)} and
$x_1=T_{\{p_1,v\}}$.  In
particular $u_3x_1\in H$, hence (with $x_1\in H$ from (iv))
$u_3\in H$.
\end{enumerate}
\end{theorem}

\begin{proof}
At each puncture, one chooses a small loop around the puncture in the
$\lambda$-sphere, lifts the family of $5$-pointed spheres
$(\Chat,\{0,\infty,1,\lambda,R(\lambda)\})$ along it, and computes the
monodromy by (Q\ref{q:pl}).  For (ii),(iii),(iv) the local equation is
\[
R(\lambda)-c=\kappa\,t^2\bigl(1+O(t)\bigr),\qquad t\coloneqq\lambda-c,\quad
\kappa=-\frac{s_i'(c)^2}{D(c)}\neq0,
\]
by \eqref{eq:squares}: the vanishing of $R-c$ is of order \emph{exactly}
$2$ — its \emph{second} derivative does not vanish — because $s_i$ has a
\emph{simple} zero at $c$ (the discriminants $5,5,-3$ are nonzero, so
$s_i'(c)\neq0$) and $D(c)\neq0$ (the roots are not in
$\{0,\tfrac12,1\}$ by Proposition~\ref{prop:R}).  As $t=e^{i\theta}$,
$0\le\theta\le2\pi$, the point $v$ therefore circles $p_i$ \emph{twice},
and the loop realizes a full double contact with
$\Delta_{\{p_i,v\}}$, giving $T_{\{p_i,v\}}^{\circ2}$.
For (v),(vi): the cycle point $\lambda\to p_{i}$ collides with contact
$1$ (the equation $\lambda=c$ itself), while $v=R(\lambda)\to
p_1=\infty$ with contact $1$ \textup{(}the poles of $R$ at $0,1$ are
simple because $D(\lambda)$ has simple zeros there while
$s_0(\lambda)^2$ does not vanish: $s_0(0)=1$, $s_0(1)=-1$\textup{)};
by (Q\ref{q:pl}) the monodromy is the product of the two commuting
twists $T_{\{p_i,p_3\}}T_{\{p_1,v\}}$.
For (i): $R$ vanishes to order exactly $2$ at
$c=\tfrac{3\pm\sqrt5}2$: writing $t=\lambda-c$ and
$\sigma(t)\coloneqq R(\lambda)/t^2$, one has
$\sigma(0)=-s_0'(c)^2/D(c)\neq0$ (by the discriminant check of
Proposition~\ref{prop:R}), so
\[
v=R(\lambda)=-s_0'(c)^2\,t^2\,\bigl(1+O(t)\bigr)\,\big/D(c) .
\]
Take a loop $\lambda(\theta)=c+\varepsilon e^{i\theta}$,
$0\le\theta\le2\pi$, with $\varepsilon>0$ so small that the disc
$|z-c|\le\varepsilon$ contains no marked point other than the limiting
$p_3=c$ and that $|v(\theta)|=|s_0'(c)|^2\varepsilon^2(1+O(\varepsilon))/|D(c)|$
is much smaller than $|c-\{0,1\}|$ and much smaller than
$\varepsilon$ — possible since the roots satisfy $c\neq0,1$ and
$v=O(\varepsilon^2)$ while $p_3-c=\varepsilon e^{i\theta}$ has modulus
$\varepsilon$.  Along this loop:
\begin{itemize}
\item $p_3(\theta)=c+\varepsilon e^{i\theta}$ winds once around $c$;
the disc bounded by its path contains \emph{no} marked point (the other
points $0,1,\infty$ are at distance $\ge\min(|c|,|1-c|)\gg\varepsilon$
and $v=O(\varepsilon^2)\ll\varepsilon$ stays off the path), so this
motion is isotopic to the identity: the $p_3$-factor contributes the
trivial mapping class.
\item $v(\theta)=\kappa\varepsilon^2 e^{2i\theta}(1+O(\varepsilon))$
winds \emph{twice} around $0=p_0$ (the argument of $v$ increases by
$4\pi$), and the circle $|z|=\kappa\varepsilon^2$ encloses only $p_0$
among the marked points ($p_3$ is at distance $\approx|c|$ from $0$).
The point-push of $v$ once around $p_0$ is the full twist
$x_0=T_{\{p_0,v\}}$ (Q\ref{q:birman}); two turns give $x_0^{\circ2}$.
\end{itemize}
Both motions are simultaneous along the single $\theta$-loop; since the
support of the $v$-push (a small annulus around $p_0$) is disjoint from
the support of the $p_3$-motion (a small annulus around $c$), the two
factors commute, and the total monodromy is
$x_0^{\circ2}\cdot\mathrm{id}=x_0^{\circ2}$.  This is the monodromy of
the simple $\lambda$-loop; hence $x_0^{\circ2}\in H$.
For (vii): set $w=\lambda^{-1}$ and take the loop
$\lambda(\theta)=\varepsilon^{-1}e^{-i\theta}$, $0\le\theta\le2\pi$
(i.e.\ $w=\varepsilon e^{i\theta}$), with $\varepsilon>0$ small.  From
the square identities \eqref{eq:squares},
\[
R(\lambda)=-\frac{\lambda^4(1+O(w))}{8\lambda^3(1+O(w))}
=-\frac{\lambda}{8}\bigl(1+O(w)\bigr)
=-\frac{e^{-i\theta}}{8\varepsilon}\bigl(1+O(\varepsilon)\bigr),
\]
so $p_3=\lambda$ and $v$ both tend to $\infty=p_1$ as
$\varepsilon\to0$, with ratio
\[
\frac{v}{p_3}\longrightarrow-\frac18
\qquad\text{and}\qquad
\frac{R(\lambda)-1}{R(\lambda)-\lambda}
=\frac{s_2(\lambda)^2}{s_3(\lambda)^2}\longrightarrow1
\quad(\text{since }s_2,s_3\text{ are monic}\ldots\text{ in fact both }
\sim\lambda^2).
\]
As $\theta$ runs from $0$ to $2\pi$: $p_3(\theta)=w^{-1}$ winds once
around $\infty=p_1$ (one full turn of $w$ around $0$), $v(\theta)$
winds once around $\infty=p_1$, and the two are mutually unlinked
(their ratio stays in a compact set avoiding $0,\infty$).  The marked
points $0$ and $1=c=p_0,p_2$ stay bounded, in fact fixed.  Choose a
small disc $\Delta$ around $\infty$ containing $p_1$; for
$\varepsilon$ small, $p_3$ and $v$ remain inside the annulus
$\{C^{-1}\varepsilon^{-1}<|z|<C\varepsilon^{-1}\}\subset\Delta$, while
$p_0,p_2$ remain outside $\Delta$.  The motion is the concatenation of:
the push of $p_3$ once around $p_1$ inside $\Delta$ — which on $B$
is the twist about the curve enclosing $\{p_1,p_3\}$, equivalently
(by taking the complementary side, $p_0,p_2$ being fixed outside
$\Delta$ together with $v$ at the moment of pushing $p_3$) the class
$u_3=T_{\{p_0,p_2\}}$, projected in $\PMCG(A)$ to
$T_{\{p_1,p_3\}}$ — and the push of $v$ once around $p_1$, which is
$x_1$.  Since the two pushes are realized on disjoint annuli
($p_3$ moving at radius $|w|^{-1}$, $v$ at radius
$|8w|^{-1}$, separated by the fixed ratio), they commute, and the
monodromy is $u_3x_1$.  Hence $u_3x_1\in H$, and with $x_1\in H$
from (iv), also $u_3=(u_3x_1)\,x_1^{-1}\in H$ and $u_3^{\circ2}\in H$.
\end{proof}

\begin{corollary}[structural form of $H$]\label{cor:Hstruct}
Let $K=\langle x_0,x_1,x_2,x_3\mid x_0x_1x_2x_3=1\rangle\cong F_3$ be
the Birman kernel \textup{(}Q\ref{q:birman}\textup{)}, free on the
\emph{consecutive triple} $\{x_0,x_1,x_2\}$ with
$x_3=(x_0x_1x_2)^{-1}$.  For $i=0,1,2$ let
$\mathrm{ch}_i:K\to\mathbb Z$ be the exponent-sum homomorphism in
$x_i$ on this free basis \textup{(}so
$\mathrm{ch}_i(x_j)=\delta_{ij}$ for $i,j\in\{0,1,2\}$, and
$\mathrm{ch}_i$ is the exponent-sum of $x_i$ in every word; it is
also the linking number of the $v$-strand with the $p_i$-strand of the
corresponding pure braid\textup{)}.  The monodromy group $H$
\textup{(}Theorem~\ref{thm:monodromy-full}\textup{)} satisfies:
\begin{enumerate}[label=\textup{(}\roman*\textup{)}]
\item the following elements lie in $H\cap K$:
\[
x_1\ (\text{census (iv)}),\qquad
x_0^{\circ2}\ (\text{(i)}),\qquad
x_2^{\circ2}\ (\text{(ii)}),\qquad
(x_0x_1x_2)^{-2}=x_3^{\circ2}\ (\text{(iii)});
\]
\item $H\to\PMCG(A)\cong F_2$ is onto \textup{(}images of
$u_0x_1,u_2x_1$ are $u_0,u_2$\textup{)} with kernel $H\cap K$;
\item $H\cap K\subseteq
\mathrm{ch}_0^{-1}(2\mathbb Z)\cap\mathrm{ch}_2^{-1}(2\mathbb Z)
\eqqcolon F$: each of the ten puncture-loop monodromies of
Theorem~\ref{thm:monodromy-full} has even $\mathrm{ch}_0$ \emph{and}
even $\mathrm{ch}_2$ — $x_1$ and
$u_0x_1,u_2x_1,u_3x_1$ have $(0,0)$; $x_0^2$ has $(2,0)$; $x_2^2$ has
$(0,2)$; $(x_0x_1x_2)^{-2}$ has $(-2,-2)$ — and
$(\mathrm{ch}_0,\mathrm{ch}_2):K\to\mathbb Z^2$ is a homomorphism.
Since every element of $H$ is a word in these loop classes
\textup{(}the loops generate $\pi_1(M_f)$\textup{)}, every element of
$H\cap K$ has even charges.  The subgroup $F$ has index $4$ in
$K=\langle x_0,x_1,x_2\rangle=F_3$;
\item $x_2\notin H$ \textup{(}$\mathrm{ch}_2(x_2)=1$ is odd, so (iii)
excludes it from $H\cap K$, and $x_2\in K$\textup{)} while
$x_2^{\circ2}\in H$.
\end{enumerate}
Consequently
\[
\bigl\langle x_1,\,x_0^{\circ2},\,x_2^{\circ2},\,(x_0x_1x_2)^{-2}
\bigr\rangle^{H}\ \le\ H\cap K\ \le\ F,
\qquad
H/(H\cap K)\xrightarrow{\ \sim\ }F_2,
\]
and $2\le[\PMCG(B):H]=[K:H\cap K]\le 4$, with the lower bound from
\textup{(}iv\textup{)}.
\end{corollary}

\begin{proof}
Items (i) and (ii) read off items (i)—(vi) of
Theorem~\ref{thm:monodromy-full} together with (Q\ref{q:birman}).
For (iii), let
\[
H_0\ \coloneqq\ \bigl\langle\,x_1,\ x_0^2,\ x_2^2,\ (x_0x_1x_2)^{-2},\
u_0,\ u_2,\ u_3\,\bigr\rangle\ \subseteq\ H:
\]
indeed the census supplies the four $K$-elements and the products
$u_i x_1$ ($i\in\{0,2\}$ from (v),(vi), $i=3$ from (vii)); since
$x_1\in H$ by (iv) of the census, these give
$u_i=(u_ix_1)\,x_1^{-1}\in H$ for $i\in\{0,2,3\}$, and conversely the
$u_i$ recover the $u_ix_1$.  Because the census loops exhaust the ten
punctures of $M_f$, $\pi_1(M_f)$ is generated by their classes (with
the single product-relation), so $H=H_0$; the displayed decomposition
is legitimate.  The projection of $H_0$ to $\PMCG(A)$ is generated
by $u_0,u_2,u_3$, hence (using
$T_{\{p_1,p_3\}}=(u_0u_2)^{-1}$ of the census item (vii)) by
$u_0,u_2$, and therefore surjects onto
$\PMCG(A)=\langle u_0,u_2\rangle$; its kernel is $H_0\cap K$.  The
seven displayed generators each have even
$(\mathrm{ch}_0,\mathrm{ch}_2)$-value as tabulated in the statement —
the $u_j$'s, whose defining curves avoid $v$, have value $(0,0)$.
Since $(\mathrm{ch}_0,\mathrm{ch}_2):K\to\mathbb Z^2$ is a
\emph{homomorphism}, its value on any word in these generators and
their inverses is the sum of the values; as every census element lies
in $H_0$, every element of $H\cap K$ arises as such a word projected
trivially, hence $H\cap K\subseteq F$.  Finally (iv) is immediate:

$\mathrm{ch}_2(x_2)=1\notin2\mathbb Z$, so $x_2\notin F\supseteq
H\cap K\ni x_2$ would contradict; more precisely
$x_2\notin H$ because were $x_2\in H$ then $x_2\in H\cap K\subseteq
F$, contradiction.  And $x_2^{\circ2}$ is a census generator.
The bounds $2\le[\PMCG(B):H]\le4$ follow from
$[K:F]=4$ and $x_2\notin H\cap K$.
\end{proof}

\begin{corollary}[index and components]\label{cor:index}
$\pi_0(\Df)$ is a nonempty subset of the finite coset space
$H\backslash\PMCG(B)$, with $2\le[\PMCG(B):H]\le 4$, hence
$1\le\#\pi_0(\Df)\le4$; in fact $\#\pi_0(\Df)\in\{2,4\}$ is even.
The twist $T_\gamma^{\circ2}=x_2^{\circ2}\in H$ acts as the identity on
$H\backslash\PMCG(B)$, hence on $\pi_0(\Df)$; the odd twist
$T_\gamma=x_2\notin H$ acts with orbits of size exactly $2$
\textup{(}no fixed component\textup{)}.
\end{corollary}

\begin{proof}
By Lemma~\ref{lem:sheets}, $\pi_0(\Df)\subseteq H\backslash\PMCG(B)$
is a nonempty, right-$H$-invariant subset (a union of full cosets).
By Corollary~\ref{cor:Hstruct}\,(iii),
$H\cap K\subseteq F$ with $[K:F]=4$, and $H$ surjects onto
$\PMCG(A)$ with kernel $H\cap K$
\textup{(}Corollary~\ref{cor:Hstruct}\,(ii)\textup{)}; hence
\[
[\PMCG(B):H]=[K:H\cap K]\le [K:F]=4,
\]
and this is $\ge2$ because $x_2\in K\setminus H$
\textup{(}Corollary~\ref{cor:Hstruct}\,(iv)\textup{)}.  The action of
$\PMCG(B)$ on $H\backslash\PMCG(B)$ is by right translation, so any
element of $H$ — in particular $x_2^{\circ2}$, which lies in $H$ by
Corollary~\ref{cor:Hstruct}\,(i) — fixes every coset, hence acts as
the identity on $H\backslash\PMCG(B)$ and on its subset $\pi_0(\Df)$.
For the odd twist: right multiplication by $x_2$ permutes the
\emph{selected} components, because $\tau_0\cdot x_2\in\Df$ — the
monodromy loop at the double zero of $s_2$ has winding number $2$, so
the two lifts of a parametrized small loop in $M_f$ starting at
$\tau_0$ and at $\tau_0\cdot x_2$ agree, whence the sheets through
$\tau_0$ and $\tau_0\cdot x_2$ are simultaneously selected
\textup{(}part \textup{(}(a)\textup{)} of
Lemma~\ref{lem:sheets}\textup{)}; and the selected set is
right-$H$-invariant, so $x_2$ takes the identity coset $H\in\pi_0(\Df)$
to $Hx_2\in\pi_0(\Df)$, and by $H$-invariance and
$x_2^2\in H$ it pairs up the whole selected set.  It has no fixed
component: $Hg=Hgx_2$ would give $gx_2g^{-1}\in H\cap K$
\textup{(}as $K$ is normal\textup{)}, but $\mathrm{ch}_2$ is
conjugacy-invariant on $\PMCG(B)$ \textup{(}pure braids fix the labels
$v,p_2$\textup{)}, so
$\mathrm{ch}_2(gx_2g^{-1})=\mathrm{ch}_2(x_2)=1$, odd, hence outside
$F\supseteq H\cap K$ \textup{(}Corollary~\ref{cor:Hstruct}\,(iii)\textup{)}.
Thus the selected components pair up under $x_2$, so
$\#\pi_0(\Df)$ is even, in $\{2,4\}$.
\end{proof}

\begin{proof}[Proof of Theorem~\ref{thm:components}]
The component count is Corollary~\ref{cor:index}; the actions there and
Proposition~\ref{prop:contenttable}\,(b)–(c) give the twist statements.
\end{proof}

\begin{proof}[Proof of Corollary~\ref{cor:answer}]
Any power of a Dehn twist preserving $\Df$ must, by
Proposition~\ref{prop:contenttable}, be an even power of $T_\gamma$;
the action of $T_\gamma^{\circ2}$ on $\pi_0(\Df)$ is the identity
\textup{(}and that of $T_\gamma$ is an involution with $2$-cycles,
on a set of at most $4$ elements\textup{)}.  Hence no infinite orbit
exists.  The
ultimate reason is the perfect-square structure of
Theorem~\ref{thm:mf}, which forces $x_2^{\circ2}=T_\gamma^{\circ2}$ to
occur as monodromy of the matched locus (Theorem~\ref{thm:monodromy-full}(ii)),
i.e.\ the locus itself knows how to loop around its own
degenerations.
\end{proof}
\end{document}
